\section{Introduction}
\label{sec:intro}

Every model in this thesis describes a population small enough that chance belongs to
the mechanism rather than to the error bars. A few individuals founding a colony, or a
few particles held inside a single compartment that may later open: in each case the
count is far from the regime in which a deterministic rate equation carries useful
information, and the questions worth asking --- will the population establish itself at
all, and if it survives, how large will it be? --- are questions about a distribution
rather than about a mean. This chapter assembles the apparatus needed to answer them.
The treatment is deliberately abstract, in counts, rates and generating functions, and
the modelling motifs that motivate it are left to the chapters that use them.

Two threads run through the material, and it is worth naming both before either is
developed.

The first is \emph{conditioning on survival}. A subcritical branching process dies out
with probability one, so its unconditional mean decays geometrically and reports almost
nothing about the realisations one would actually observe. Restrict attention to those
realisations still alive at generation $n$, however, and a different picture appears:
the conditional population settles onto a fixed distribution --- the
\emph{quasi-stationary} distribution --- whose mean is a constant. For the binary
Galton--Watson process with replication probability $p<\tfrac12$ that constant is
$1/A(p)$, where
\begin{equation*}
A(p) = \lim_{n\to\infty}\frac{S_n}{(2p)^n}
\end{equation*}
and $S_n$ is the probability of survival to generation $n$. A large part of this
chapter is an account of $A(p)$: what it is, how to compute it, and why, despite a
considerable search, no elementary formula for it has been found.

The second thread is \emph{the near-critical regime}. As $p$ rises towards $\tfrac12$
the process takes ever longer to die, the quasi-stationary mean $1/A(p)$ diverges, and
any numerical scheme that iterates to stationarity becomes impractical. Just above
$\tfrac12$ the size of the founding cohort begins to weigh heavily on whether the
process survives at all. The behaviour on either side of $p=\tfrac12$ is what makes
small founding numbers interesting, and it is also what makes them awkward to compute
with.

\Cref{sec:tools} sets out the standard machinery the rest of the chapter uses:
continuous-time Markov chains and their generators, exponential clocks and the
competition between them, generating functions, and the linear birth--death process
solved in closed form. Readers already fluent in that material may read it as a fixing
of notation. \Cref{sec:gw} introduces the Galton--Watson process and establishes the
two facts on which everything later rests --- that extinction is certain at and below
criticality, and that at criticality the expected lifetime is nonetheless infinite.
\Cref{sec:qsd} conditions on survival, derives the limiting conditional mean and
variance, and so introduces $A(p)$. \Cref{sec:small} asks what changes when the process
starts from a cohort of $k$ individuals rather than one, records a discrete
birth--death--catastrophe calculation of the kind developed properly in later chapters,
and works out the corresponding constant for the continuous-time birth--death process,
where a closed form \emph{is} available. \Cref{sec:Ap} is devoted to $A(p)$ itself: an
infinite-product representation, an exact series with two-sided bounds that pin the
constant between the continuous-time answer and twice it, the near-critical asymptotic
$A(p)\sim2(1-2p)$, and an account of why the exact function resists closed-form
description --- a question that turns out to concern the Koenigs linearisation of the
logistic map. \Cref{sec:moc} then changes tack and develops generating-function partial
differential equations and the method of characteristics for the absorption models used
later in the thesis.

Standard results are presented as standard. Where an argument is heuristic, or where
the numerics are doing the real work, the text says so at that point rather than in a
general disclaimer, and every claim that can be checked numerically has been.
